\section{Trace anomaly}\label{sec:traceanom}

As a test of our result, we use the trace anomaly of the \emt. As
suggested in \citere{Nielsen:1977sy}, a simple derivation consists of
taking the trace of the \emt\ in $D=4 - 2 \ep$ dimensions. By use of the
equations of motion, this gives for the gauge invariant part
\begin{equation}
	T_{\mu \mu} = \frac{\ep}{2 g_0^2} F_{\rho \sigma}^a F_{\rho
          \sigma}^a - \sum_{f=1}^{n_F} m_{f, 0} \overline{\psi}_f \psi_f =
        \frac{1}{2D}\left(\frac{\ep}{g_0^2}
        \calo_{2,\mu\mu}+\calo_{4,\mu\mu}\right)
        \,,
	\label{eq:EMTabare}
\end{equation}
where we have used \eqn{eq:eom} in the last step.
Using
the mixing matrix $\zeta_{i j}(t)$, we can rewrite this in terms of
flowed operators:
\begin{equation}
\begin{split}
  T_{\mu\mu} =  \bar c_i(t)\tcalo_{i,\mu\mu}(t,x)\,,\qquad
  \bar c_i(t) = \frac{1}{2D}\left(
  \frac{\ep}{ g_0^2} \zeta_{2 i}^{-1}(t)
         + \zeta_{4 i}^{-1}(t)\right)\,.
	 \label{eq:EMTaflowed}
  \end{split}
\end{equation}
Note that $\bar{c}_1 (t) = \bar{c}_3 (t) = 0$, as $\tcalo_{1, \mu \nu} (t, x)$
and $\tcalo_{3, \mu \nu} (t, x)$ have a non-trivial index structure and
therefore $\calo_{2, \mu \nu} (x)$ and $\calo_{4, \mu \nu} (x)$ cannot
mix with them.
Since $\tcalo_{2,\mu\mu}=D\tcalo_{1,\mu\mu}$ and
$2\tcalo_{4,\mu\mu}=D\tcalo_{3,\mu\mu}$, we cannot equate coefficients
with \eqn{eq:emtflow} for all $i$ individually. Instead, only the weaker
conditions
\begin{equation}
  \begin{split}
    c_1(t)+Dc_2(t) = D\bar c_2(t)\,,\qquad
    2c_3(t)+Dc_4(t) = D\bar c_4(t)\,.
  \end{split}
\end{equation}
can be derived. We checked that these equations are indeed fulfilled by
our result.
